\documentclass[12pt,a4paper]{article}

\usepackage[utf8]{inputenc}
\usepackage[T1]{fontenc}
\usepackage[brazil]{babel}
\usepackage{lmodern}
\usepackage{geometry}
\usepackage{setspace}
\usepackage{longtable}
\usepackage{tabularx}
\usepackage{array}
\usepackage{booktabs}
\usepackage{enumitem}
\usepackage{hyperref}
\usepackage{titlesec}

\geometry{margin=2.5cm}
\setstretch{1.15}
\setlength{\parindent}{0pt}
\setlength{\parskip}{0.6em}
\hypersetup{
  colorlinks=true,
  linkcolor=black,
  urlcolor=blue,
  pdftitle={Proposta Tecnica e Comercial com Memorial Descritivo},
  pdfauthor={OpenAI Codex}
}

\titleformat{\section}{\normalfont\bfseries\uppercase}{\thesection.}{0.5em}{}
\titleformat{\subsection}{\normalfont\bfseries}{\thesubsection}{0.5em}{}

\newcolumntype{Y}{>{\raggedright\arraybackslash}X}

\begin{document}

\begin{center}
{\Large \textbf{PROPOSTA TÉCNICA E COMERCIAL COM MEMORIAL DESCRITIVO}}\\[1em]
\textbf{NOME DO SISTEMA: Nobly, SIMAS PRO}
\end{center}

\textbf{Objeto de referência:} licenciamento, implantação, parametrização, treinamento, suporte técnico e manutenção do sistema para atendimento especializado do SUAS, no âmbito do CREAS e da Vigilância Socioassistencial do Município de Patos/PB, com foco em PAEFI, MSE e inteligência socioassistencial.

\textbf{Delimitação de escopo:} esta proposta considera exclusivamente os módulos e funcionalidades relacionados ao CREAS, às Medidas Socioeducativas em Meio Aberto e à Vigilância Socioassistencial, desconsiderando os demais módulos existentes na plataforma.

\section{OBJETO}

A presente proposta tem por objeto a disponibilização de licença de uso do sistema \textbf{Nobly, SIMAS PRO}, já desenvolvido, funcional e vocacionado às demandas da política pública de Assistência Social, para utilização no Município de Patos/PB, especificamente nos serviços, fluxos e rotinas do \textbf{CREAS} e da \textbf{Vigilância Socioassistencial}.

O objeto compreende, de forma integrada:

\begin{itemize}[leftmargin=1.5cm]
\item a \textbf{licença de uso} do sistema em ambiente web, para operação institucional pelas equipes autorizadas;
\item a \textbf{implantação} da solução, com preparação do ambiente, habilitação dos módulos contratados e adequação operacional ao fluxo local;
\item a \textbf{parametrização} dos recursos necessários ao uso do sistema nas rotinas do CREAS e da Vigilância Socioassistencial;
\item o \textbf{treinamento} dos usuários e da equipe gestora, com orientação prática voltada ao uso diário da ferramenta;
\item o \textbf{suporte técnico} continuado, para atendimento de demandas operacionais, correção de inconsistências e sustentação do uso;
\item a \textbf{manutenção evolutiva}, destinada à preservação da aderência funcional da solução ao contexto de uso especializado.
\end{itemize}

No âmbito material desta contratação, a solução destina-se a:

\begin{itemize}[leftmargin=1.5cm]
\item operacionalizar o prontuário digital do atendimento especializado no CREAS, com foco em PAEFI;
\item apoiar o acompanhamento de \textbf{Medidas Socioeducativas em Meio Aberto}, com controle individualizado dos registros;
\item viabilizar a leitura analítica da demanda socioassistencial por meio do \textbf{Painel de Vigilância Socioassistencial}, com indicadores, mapas, gráficos e drill-down operacional.
\end{itemize}

\section{JUSTIFICATIVA DE SINGULARIDADE TÉCNICA}

O sistema \textbf{Nobly, SIMAS PRO} apresenta singularidade técnica no contexto desta contratação por se tratar de solução já estruturada para o \textbf{SUAS}, com desenho funcional aderente à rotina de atendimento especializado, monitoramento de casos e produção de inteligência socioassistencial. A documentação técnica interna do projeto identifica explicitamente a solução como um \textbf{sistema SUAS Patos/PB}, evidenciando alinhamento concreto com o contexto municipal e com a organização dos fluxos locais.

Sua especialização não decorre apenas da existência de telas ou formulários, mas da forma como o sistema articula, em uma mesma base operacional, o cadastro qualificado do caso, o prontuário digital, os encaminhamentos, os registros de acompanhamento, os dados familiares, os marcadores de vulnerabilidade, os registros de MSE e a camada analítica de Vigilância. Essa integração reduz retrabalho, elimina dispersão de informação e amplia a capacidade de coordenação entre atendimento técnico, gestão e monitoramento territorial.

No âmbito do CREAS, a solução foi desenhada para refletir a lógica de prontuário socioassistencial e de acompanhamento contínuo, com registro de histórico, composição familiar, vínculos, agravos, encaminhamentos e evolução do caso. No âmbito da Vigilância Socioassistencial, o sistema consolida indicadores territoriais e operacionais em painéis dinâmicos, permitindo leitura gerencial e apoio à tomada de decisão com base na própria base viva do atendimento.

Há, ainda, forte dificuldade de substituição por sistemas genéricos. Soluções inespecíficas, ainda que tecnicamente robustas em sentido amplo, normalmente não contemplam a modelagem do dado socioassistencial com a profundidade necessária para PAEFI, MSE e análise territorial. No presente caso, a singularidade técnica reside justamente na combinação entre:

\begin{itemize}[leftmargin=1.5cm]
\item aderência ao domínio do SUAS;
\item aderência ao contexto institucional de Patos/PB;
\item integração entre operação e inteligência;
\item estrutura analítica orientada ao CREAS e à Vigilância;
\item baixa substituibilidade por plataformas genéricas sem perda funcional, metodológica e informacional.
\end{itemize}

\section{FUNDAMENTAÇÃO JURÍDICA E VALORES}

A presente contratação encontra fundamento no \textbf{art. 75, inciso II, da Lei nº 14.133/2021}, que admite a dispensa de licitação para contratação que envolva valores inferiores ao limite legalmente atualizado para outros serviços e compras.

Para fins de enquadramento jurídico, registra-se que o limite atualizado do \textbf{art. 75, inciso II}, vigente a partir de \textbf{1º de janeiro de 2026}, corresponde a \textbf{R\$ 65.492,11}, conforme atualização oficial promovida pelo \textbf{Decreto nº 12.807, de 29 de dezembro de 2025}. Assim, o valor global desta proposta, de \textbf{R\$ 65.400,00}, permanece dentro do teto legal aplicável à contratação direta por dispensa, sem extrapolação do limite normativo.

\subsection*{PARTE A --- IMPLANTAÇÃO, PARAMETRIZAÇÃO E TREINAMENTO}

\begin{center}
\begin{tabularx}{\textwidth}{Y>{\raggedleft\arraybackslash}p{2.8cm}>{\raggedleft\arraybackslash}p{3.2cm}}
\toprule
\textbf{Etapa} & \textbf{Percentual} & \textbf{Valor} \\
\midrule  
implantação & 40\% & R\$ 10.400,00 \\
Parametrização & 35\% & R\$ 9.100,00 \\
Treinamento & 25\% & R\$ 6.500,00 \\
\textbf{Total Parte A} &  & \textbf{R\$ 26.000,00} \\
\bottomrule
\end{tabularx}
\end{center}

\subsection*{PARTE B --- SUPORTE, MANUTENÇÃO E LICENÇA}

\begin{center}
\begin{tabularx}{\textwidth}{Y>{\raggedleft\arraybackslash}p{4cm}}
\toprule
\textbf{Item} & \textbf{Valor} \\
\midrule
Licença de uso & R\$ 1.400,00 \\
Suporte técnico & R\$ 900,00 \\
Manutenção evolutiva & R\$ 1.000,00 \\
\textbf{Total mensal} & \textbf{R\$ 3.300,00} \\
\textbf{Total anual (12 meses)} & \textbf{R\$ 39.600,00} \\
\bottomrule
\end{tabularx}
\end{center}

\subsection*{RESUMO FINAL}

\begin{center}
\begin{tabularx}{\textwidth}{Y>{\raggedleft\arraybackslash}p{4cm}}
\toprule
\textbf{Composição} & \textbf{Valor} \\
\midrule
Parte A & R\$ 26.000,00 \\
Parte B & R\$ 39.600,00 \\
\textbf{VALOR TOTAL} & \textbf{R\$ 65.400,00} \\
\bottomrule
\end{tabularx}
\end{center}

Sob o ponto de vista técnico, o valor proposto corresponde a uma solução já funcional, especializada no SUAS, com capacidade de elevar a eficiência operacional do CREAS e da Vigilância Socioassistencial, qualificar a produção de dados, apoiar a inteligência territorial e reduzir perdas decorrentes de controles manuais ou de sistemas genéricos. Trata-se, portanto, de solução com valor agregado intrinsecamente vinculado à sua aderência temática e funcional, não sendo tecnicamente comparável a plataformas genéricas sem especialização socioassistencial.

\section{MÓDULOS E FUNCIONALIDADES DO SISTEMA}

\subsection{Gestão de Prontuários Digitais (PAEFI/CREAS)}

O módulo de prontuário digital do CREAS foi estruturado para o registro qualificado do atendimento especializado, servindo como núcleo operacional do caso. A documentação técnica do sistema demonstra fluxo próprio de criação, edição e detalhamento do prontuário, permitindo a consolidação do histórico do atendimento e sua continuidade técnica.

Entre as funcionalidades centrais, destacam-se:

\begin{itemize}[leftmargin=1.5cm]
\item registro completo dos atendimentos e dados estruturantes do caso;
\item consolidação do histórico familiar e dos elementos de composição do núcleo;
\item acompanhamento da evolução do caso ao longo do tempo;
\item registro de encaminhamentos, confirmações, reincidências e desdobramentos do atendimento;
\item manutenção do prontuário em tela de detalhe, com continuidade operacional do caso.
\end{itemize}

Quanto à profundidade informacional, a arquitetura documental do módulo mapeia \textbf{51 campos de prontuário e formulário}, abrangendo campos de topo, dados de atendimento, vítima, família, saúde, encaminhamentos, agressor e moradia, além de campos condicionais acionados conforme a situação concreta. Na prática, isso significa o registro de \textbf{dezenas de variáveis socioassistenciais}, incluindo dados familiares, dados territoriais, marcadores de vulnerabilidade, benefícios, confirmação da violência, reincidência, composição do domicílio, condição de saúde, vínculo com agressor, moradia e demais elementos necessários ao acompanhamento técnico qualificado.

Essa riqueza de dados qualifica o PAEFI não apenas como fluxo de atendimento, mas como base estruturada de conhecimento sobre a trajetória do caso, permitindo leitura individual, familiar e territorial a partir de uma mesma matriz informacional.

\subsection{Módulo de Medidas Socioeducativas (MSE)}

O módulo de MSE é destinado ao controle especializado das Medidas Socioeducativas em Meio Aberto, com foco no acompanhamento individualizado do adolescente e no monitoramento do cumprimento da medida.

As funcionalidades identificadas na documentação e na estrutura do módulo compreendem:

\begin{itemize}[leftmargin=1.5cm]
\item cadastro e listagem de registros de medidas socioeducativas;
\item controle do tipo de medida aplicada;
\item acompanhamento individualizado do adolescente e de seu responsável;
\item registro da situação de \textbf{cumprimento} ou \textbf{descumprimento};
\item indicação do local de descumprimento, quando aplicável;
\item registro do PIA, com data de elaboração e status;
\item visualização em modal de detalhe e edição.
\end{itemize}

Cada registro de MSE monitora, no mínimo, \textbf{13 informações estruturadas por caso}, tais como nome do adolescente, data de nascimento, responsável legal, endereço, contato, NIS, tipo de medida, data de início, duração em meses, situação atual, local de descumprimento, data do PIA e status do PIA. Além disso, o próprio módulo calcula e exibe indicadores derivados relevantes ao acompanhamento, como \textbf{idade atual} e \textbf{data estimada de término} da medida, compondo uma linha do tempo operacional do adolescente.

No plano gerencial, o módulo disponibiliza \textbf{4 indicadores de acompanhamento}, voltados a total de medidas ativas, casos em cumprimento, casos em descumprimento e medidas com vencimento próximo. Isso favorece a identificação de evolução, criticidade, risco de descumprimento e necessidade de intervenção tempestiva, inclusive para leitura de reincidência e recorrência de fragilidades no acompanhamento.

\subsection{Painel de Vigilância Socioassistencial}

O Painel de Vigilância Socioassistencial constitui a camada analítica da solução, transformando a base operacional do atendimento em informação gerencial e territorialmente orientada.

Conforme a arquitetura do sistema, o painel agrega, de forma paralela e dinâmica, \textbf{6 blocos analíticos centrais}:

\begin{itemize}[leftmargin=1.5cm]
\item fluxo de demanda;
\item sobrecarga da equipe;
\item incidência por bairros;
\item fontes de acionamento;
\item reincidência;
\item perfil das violações.
\end{itemize}

Do ponto de vista visual e funcional, o painel opera com \textbf{múltiplos painéis dinâmicos}, utilizando mapa territorial, gráficos de barras, gráficos de pizza e modal de drill-down para abertura da lista de casos correspondente ao insight selecionado. A documentação do backend do dashboard também demonstra capacidade analítica composta por \textbf{8 indicadores principais}, \textbf{4 indicadores de contexto familiar} e \textbf{7 gráficos gerenciais}, o que reforça a robustez da camada de análise que alimenta o uso institucional da informação.

Em termos de atualização, a solução trabalha sobre a base operacional do sistema, permitindo leitura em tempo real do acervo cadastrado e dos movimentos registrados no atendimento. Assim, a cada inserção, atualização ou consolidação de caso, o painel pode refletir o comportamento da demanda, a concentração territorial, os perfis de violação, os canais de acionamento e a reincidência, oferecendo suporte efetivo à gestão, ao planejamento e à priorização de ações.

A capacidade analítica do módulo é especialmente relevante para a Administração por permitir:

\begin{itemize}[leftmargin=1.5cm]
\item leitura territorial da incidência socioassistencial;
\item acompanhamento da pressão de demanda sobre as equipes;
\item identificação de padrões de violação e recorrência;
\item apoio à tomada de decisão baseada em evidências;
\item integração entre visão operacional e visão gerencial.
\end{itemize}

\section{VANTAGENS PARA A ADMINISTRAÇÃO}

\subsection*{Eficiência administrativa e operacional}

A solução reduz retrabalho, centraliza informações e organiza em fluxo único o registro, a consulta e o monitoramento de casos, fortalecendo a produtividade das equipes do CREAS e o acompanhamento gerencial da Vigilância Socioassistencial.

\subsection*{Redução de custos indiretos}

Ao substituir controles dispersos, registros manuais e consolidações paralelas de dados, o sistema diminui o custo administrativo associado a retrabalho, perda de informação, consolidação tardia de relatórios e baixa rastreabilidade dos atendimentos.

\subsection*{Segurança de dados e segregação de acesso}

A arquitetura documentada contempla autenticação, rotas protegidas, controle granular de permissões e recorte de acesso por unidade, o que favorece tratamento mais seguro das informações socioassistenciais e uso institucional com perfis diferenciados.

\subsection*{Transparência, rastreabilidade e gestão}

O prontuário digital, os registros estruturados e a camada analítica ampliam a rastreabilidade do atendimento, facilitam supervisão técnica, produção de relatórios e tomada de decisão baseada em evidências extraídas da própria base municipal.

\subsection*{Melhoria do atendimento ao usuário do SUAS}

Com dados mais completos, histórico acessível, organização dos encaminhamentos e monitoramento territorial, a Administração amplia a capacidade de resposta qualificada às situações acompanhadas, melhorando a continuidade do atendimento e a priorização de casos.

\section{CONCLUSÃO}

À vista do exposto, conclui-se que a contratação do sistema \textbf{Nobly, SIMAS PRO}, no escopo restrito ao \textbf{CREAS} e à \textbf{Vigilância Socioassistencial}, mostra-se tecnicamente vantajosa, juridicamente enquadrável e alinhada ao interesse público municipal.

Sob o aspecto jurídico, o valor global de \textbf{R\$ 65.400,00} enquadra-se no limite de \textbf{R\$ 65.492,11} aplicável ao \textbf{art. 75, inciso II, da Lei nº 14.133/2021}, legitimando a adoção da dispensa de licitação, desde que observadas as formalidades próprias do processo administrativo de contratação direta.

Sob o aspecto técnico, a solução apresenta aderência específica ao SUAS, à rotina do CREAS, ao controle de MSE e à produção de inteligência socioassistencial, com nível de especialização e integração que supera soluções genéricas. Sob o aspecto material, atende ao interesse público por fortalecer a capacidade institucional do Município de Patos/PB para registrar, acompanhar, monitorar e analisar sua demanda socioassistencial de forma qualificada, segura e tempestiva.

\section{MEMORIAL DESCRITIVO}

\subsection{Arquitetura do sistema}

O sistema possui arquitetura de aplicação web em modelo monolítico com separação lógica entre frontend e backend. A camada de apresentação opera como \textbf{SPA (Single Page Application)}, responsável pela navegação protegida, formulários, painéis analíticos e interação do usuário. A camada de backend atua como \textbf{API} de serviços, concentrando autenticação, regras de acesso, lógica de negócio e consulta aos dados.

Essa arquitetura é adequada ao objeto porque permite:

\begin{itemize}[leftmargin=1.5cm]
\item operação centralizada em ambiente web;
\item atualização contínua sem reinstalação local;
\item compartilhamento de base única entre operação e análise;
\item integração entre prontuário, MSE e Vigilância Socioassistencial.
\end{itemize}

\subsection{Módulos contemplados}

No escopo desta proposta, o sistema compreende os seguintes módulos especializados:

\begin{itemize}[leftmargin=1.5cm]
\item \textbf{Prontuário Digital CREAS/PAEFI}, para cadastro, atualização, histórico, encaminhamentos e detalhamento do caso;
\item \textbf{Controle de MSE}, para registro, busca, monitoramento e acompanhamento das medidas socioeducativas em meio aberto;
\item \textbf{Painel de Vigilância Socioassistencial}, para monitoramento gerencial e territorial, com indicadores, mapas, gráficos e drill-down.
\end{itemize}

\subsection{Tecnologias}

De acordo com a documentação técnica do projeto, a solução utiliza tecnologias modernas e consolidadas para aplicações web institucionais, destacando-se:

\begin{itemize}[leftmargin=1.5cm]
\item frontend em aplicação web responsiva;
\item backend em API de serviços;
\item banco de dados relacional com suporte a estruturas semiestruturadas para o prontuário;
\item mecanismos de validação de formulário e adaptação de contratos de dados;
\item bibliotecas de visualização analítica para gráficos e mapa territorial.
\end{itemize}

Na base documentada do sistema, identificam-se, entre outras, as tecnologias \textbf{React}, \textbf{TypeScript}, \textbf{Vite}, \textbf{Express} e \textbf{PostgreSQL}, compondo uma pilha adequada para sustentação, evolução e operação institucional.

\subsection{Segurança}

No aspecto de segurança e controle de acesso, o sistema conta com:

\begin{itemize}[leftmargin=1.5cm]
\item autenticação por token;
\item rotas protegidas no frontend;
\item controle granular de permissões por tela e por operação;
\item recorte de acesso por unidade;
\item tratamento de registros com exclusão lógica;
\item segregação entre usuários gestores e usuários vinculados à unidade.
\end{itemize}

Tais elementos contribuem para a proteção das informações sensíveis tratadas pelo CREAS e para a governança do uso dos dados na Vigilância Socioassistencial.

\subsection{Escalabilidade}

A solução apresenta escalabilidade funcional e evolutiva, pois sua arquitetura separa camada de interface, serviços e persistência de dados, permitindo expansão de campos, filtros, indicadores, relatórios e rotinas de monitoramento sem ruptura do núcleo operacional.

Também se verifica capacidade de crescimento no uso institucional, uma vez que:

\begin{itemize}[leftmargin=1.5cm]
\item a base do prontuário admite ampliação de atributos socioassistenciais;
\item a camada analítica admite incorporação de novos indicadores;
\item o modelo de acesso por unidade favorece organização institucional do uso;
\item os módulos compartilham contratos e serviços reutilizáveis, reduzindo impacto de evolução futura.
\end{itemize}

\section*{Base documental utilizada}

Esta proposta foi elaborada com base nos documentos técnicos existentes no repositório do sistema, especialmente:

\begin{itemize}[leftmargin=1.5cm]
\item \texttt{resumo.md};
\item \texttt{frontend/docs/arquitetura\_front.md};
\item \texttt{frontend/docs/cadastro/arquitetura\_cadastro.md};
\item \texttt{frontend/docs/cadastro/cadastro\_models.md};
\item \texttt{frontend/docs/cadastro/formulario-cadastro.md};
\item \texttt{frontend/docs/cadastro/inventario\_cadastro\_legacy.md};
\item \texttt{frontend/docs/refactor/dashboard-drilldown-contrato-atual.md};
\item \texttt{frontend/docs/migrar\_dashboard\_vigilancia.md};
\item \texttt{frontend/docs/cadastro/frontend-migracao-drilldown-vigilancia.md};
\item \texttt{frontend/docs/refactor/refactor\_graficos.md}.
\end{itemize}

\section*{Referências normativas oficiais}

\begin{itemize}[leftmargin=1.5cm]
\item Lei nº 14.133/2021: \url{https://www.planalto.gov.br/ccivil_03/_ato2019-2022/2021/lei/l14133.htm}
\item Decreto nº 12.807/2025: \url{https://www.planalto.gov.br/ccivil_03/_ato2023-2026/2025/decreto/d12807.htm}
\end{itemize}

\end{document}
